\documentclass[11pt]{article}
\usepackage{mathblatt}
% Prüfungsblatt nach mathe-nachhilfe pruefung.md § 5 (Runde Dreiecke, 11.10.2026), Hand-Satz; Präambel des Setzers; Muster P4Z.
% Präambel des Setzers (werkzeuge/setzer.py), wortgleich aus dem Hand-Satz T6B (bau/proben/2026-10-09/kathete-berechnen), dort aus K4W.
\makeatletter
\setlength{\mbpfbe}{0mm}% keine Punkte (Bauregel 6.4)
% eingerückt (Bauregel 4.7, 6.6): \gEin Einzug, \gSchrift eine Stufe kleiner
\newlength{\gEin}\setlength{\gEin}{0pt}
\let\gSchrift\relax
\renewcommand{\pfkreuzzeile}[1]{\par\noindent\hangindent3.2em\hangafter1\hspace*{1.6em}$\square$\ #1\par}
\newcommand{\gkreuz}[1]{$\square$\ #1\hspace{2em}}
% Bauregel 6.7: links, was man ansieht, rechts, was man tut; Spaltengrenze überall an derselben Stelle
\newsavebox{\gbox}\newlength{\gLinks}
\newcommand{\gzwei}[2]{\par\smallskip\noindent\sbox{\gbox}{#1}%
  \setlength{\gLinks}{\dimexpr0.5\textwidth-\gEin-\mbpfnr\relax}%
  \ifdim\wd\gbox>\dimexpr\gLinks-4mm\relax
    \usebox{\gbox}\par\smallskip\noindent\begin{minipage}{\linewidth}#2\end{minipage}%
  \else
    \begin{minipage}[t]{\gLinks}\vspace{0pt}\usebox{\gbox}\end{minipage}%
    \begin{minipage}[t]{\dimexpr\linewidth-\gLinks\relax}\vspace{0pt}\raggedright #2\end{minipage}%
  \fi\par}
\RenewDocumentEnvironment{pfaufg}{m m m}{%
  \begin{lrbox}{\mbaufgabebox}\begin{minipage}[t]{\linewidth}%
  \gSchrift\noindent\hspace*{\gEin}\mbpfrandmarke{#1}%
  \begin{minipage}[t]{\mbpfnr}\strut\textbf{#2}\end{minipage}%
  \begin{minipage}[t]{\dimexpr\linewidth-\gEin-\mbpfnr\relax}\raggedright\strut\ignorespaces}%
 {\end{minipage}%
  \end{minipage}\end{lrbox}%
  \setlength{\mbaufgabehoehe}{\dimexpr\ht\mbaufgabebox+\dp\mbaufgabebox\relax}%
  \par\addvspace{7pt}\Needspace*{\mbaufgabehoehe}%
  \noindent\usebox{\mbaufgabebox}\par\addvspace{7pt}}
\makeatother
% Rechenraum als Karo (Bauregel 6.9): n Rechenschritte = n mal 10 mm
\newcommand{\karo}[1]{\par\smallskip\noindent\begin{tikzpicture}
  \clip (0,0) rectangle ({\linewidth-1mm},{#1*10mm});
  \draw[black!16,line width=0.3pt,step=5mm] (0,0) grid ({\linewidth},{#1*10mm});
  \end{tikzpicture}\par}
\newcommand{\kk}[1]{$\square$\,#1\hspace{0.9em}}
\newcommand{\linie}[1]{\,{\color{black!45}\rule[-1pt]{#1}{0.4pt}}\,}
% Zeile am Ende jedes Durchgangs: Originale klein grau, darüber; dann Vorgänger/Nachfolger (Bauregel 3.4, 6.5)
\newcommand{\blattende}[2]{\par\vfill\noindent\begin{minipage}{\linewidth}{\scriptsize\color{mbgrau}Originale: #1\par}\smallskip
{\footnotesize #2\par}\end{minipage}}
\newcommand{\pkt}{\textperiodcentered{}}

\newcommand{\kennung}{T2U}
% Teilaufgabe mit Skizze links in gzwei: Buchstabe in eigener Spalte, darunter Text und Skizze
\makeatletter
\newcommand{\gteil}[2]{\makebox[\mbpfteil][l]{\textbf{#1}}\parbox[t]{\dimexpr0.5\textwidth-\gEin-\mbpfnr-5mm-\mbpfteil\relax}{\raggedright #2}}
\makeatother

\begin{document}
\pfheftstil
\pflaufkopf{Winkel}{\kennung}\pfseitevon
\pfheftkopf[\kennung]{Winkel}{P10}
\pfstufenzeile{Eigenschaften kennen \pkt{} Innenwinkelsumme: Dreieck \pkt{} Viereck \pkt{} gleichschenklig: zwei gleiche Winkel \pkt{} rechter Winkel: begründen}

% ===================== Stufe 1 =====================
\Needspace*{0.25\textheight}\pfstufekopf{Eigenschaften kennen}{letzte 5 Jahre: 2026}

% Bank winkel-dreiecke-e3-k2-s10-v2 (Raute; Verfremdung von 2016-OS-B1e)
\begin{pfaufg}{}{1.}{}
Kreuze an, welche Aussage für jede Raute gilt.
\gzwei{\begin{tikzpicture}[line width=0.7pt]\draw (0,0)--(1.6,0)--(2.6,1.249)--(1.0,1.249)--cycle;\end{tikzpicture}}{\pfkreuzab{A}{Benachbarte Winkel sind gleich groß.}
\pfkreuzab{B}{Gegenüberliegende Winkel sind gleich groß.}
\pfkreuzab{C}{Alle vier Winkel sind gleich groß.}}
\fusshilfe{\mbox{1:~B}}
\end{pfaufg}

% 2026-FOR-B1c (EBR-Zwilling 2026-EBR-B1c), eigener Wortlaut
\begin{pfaufg}{nach~P10\\’26~\pkt{}~1c}{2.}{}
Kreuze an, welche Aussage für jedes Trapez gilt.\par\smallskip
\pfkreuzab{A}{Alle vier Winkel sind gleich groß.}
\pfkreuzab{B}{Mindestens zwei Seiten sind zueinander parallel.}
\pfkreuzab{C}{Alle vier Seiten sind gleich lang.}
\fusshilfe{\mbox{2:~B}}
\end{pfaufg}

% ===================== Stufe 2 =====================
\Needspace*{0.25\textheight}\pfstufekopf{Innenwinkelsumme: Dreieck \pkt{} Viereck}{letzte 5 Jahre: 2026 \pkt{} 2023 \pkt{} 2022}

% Bank winkel-dreiecke-e3-k2-s1-v1, e3-k2-s6-v1, c) eigen (Päckchen: Skizze → Text → krumm, Durchsicht 11.10.)
\begin{pfaufg}{}{3.}{}
Berechne den fehlenden Winkel.\par\smallskip
\gzwei{\gteil{a)}{Dreieck\par\smallskip\begin{tikzpicture}[line width=0.7pt]\draw (0,0)--(2.6,0)--(1.805,2.310)--cycle;\draw[mbgrau,line width=0.5pt] (0.4,0) arc (0:52:0.4);\node[inner sep=0pt,font=\footnotesize] at (0.66,0.31) {$52^\circ$};\draw[mbgrau,line width=0.5pt] (2.2,0) arc (180:109:0.4);\node[inner sep=0pt,font=\footnotesize] at (2.0,0.36) {$71^\circ$};\end{tikzpicture}}}{\karo{2}}
\gzwei{\gteil{b)}{Viereck mit den Winkeln $84^\circ$, $97^\circ$ und $103^\circ$}}{\karo{2}}
\gzwei{\gteil{c)}{Viereck mit den Winkeln $86{,}5^\circ$, $97{,}3^\circ$ und $101{,}4^\circ$}}{\karo{2}}
\fusshilfe{\mbox{3:~a) 57°; b) 76°; c) 74,8°}}
\end{pfaufg}

% Bank winkel-dreiecke-e3-k2-s7-v7 (Teildreieck an der Höhe), Durchsicht 11.10.: ohne Skizze, Text
\begin{pfaufg}{}{4.}{}
Im Dreieck $ABC$ ist $\alpha = 52^\circ$. Die Höhe $h_c$ geht von $C$ senkrecht auf die Seite $\overline{AB}$. Berechne den Winkel $\varepsilon$ zwischen $\overline{AC}$ und der Höhe.
\karo{2}
\fusshilfe{\mbox{4:~38°}}
\end{pfaufg}

% 2026-FOR-B1i (EBR-Zwilling 2026-EBR-B1i), eigener Wortlaut
\begin{pfaufg}{nach~P10\\’26~\pkt{}~1i}{5.}{}
Die Figur zeigt das Parallelogramm $ABCD$. Ein Winkel ist $75^\circ$ groß. Bestimme den Winkel $\beta$.
\gzwei{\begin{tikzpicture}[line width=0.7pt]\draw (0, 0)--(3.2, 0)--(3.614,1.545)--(0.414,1.545)--cycle;\draw[mbgrau,line width=0.5pt] (0.380,0.000) arc (0.00:75.00:0.38);\node[inner sep=0pt,font=\footnotesize] at (0.571,0.438) {$75^\circ$};\draw[mbgrau,line width=0.5pt] (3.298,0.367) arc (75.00:180.00:0.38);\node[inner sep=0pt,font=\footnotesize] at (2.804,0.516) {$\beta$};\node[font=\small,anchor=north east] at (0.000,0.000) {$A$};\node[font=\small,anchor=north west] at (3.200,0.000) {$B$};\node[font=\small,anchor=south west] at (3.614,1.545) {$C$};\node[font=\small,anchor=south east] at (0.414,1.545) {$D$};\end{tikzpicture}}{\karo{2}}
\fusshilfe{\mbox{5:~105°}}
\end{pfaufg}

% 2022-OS-K5c, herausgelöst (Seite b, Höhe 7,4 m und 52° bei B nur für 5a, 5b, 5d, weggelassen)
\begin{pfaufg}{nach~P10\\’22~\pkt{}~5c}{6.}{}
Das Dreieck $ABC$ hat keinen rechten Winkel. Die Höhe $h_c$ geht von $C$ senkrecht auf die Seite $\overline{AB}$, ihr Fußpunkt heißt $D$. Der Winkel $\alpha$ ist etwa $31{,}7^\circ$ groß. Ermittle die Größe des Winkels $\gamma_1$ zwischen $\overline{AC}$ und der Höhe.
\gzwei{\begin{tikzpicture}[line width=0.7pt]\draw (0, 0)--(4.45, 0)--(3.0, 1.85)--cycle;\draw[dashed] (3.0, 1.85)--(3.0, 0);\draw[mbgrau,line width=0.5pt] (3.300,0.000) arc (0.00:90.00:0.3);\fill (3.092,0.092) circle (0.8pt);\draw[mbgrau,line width=0.5pt] (0.600,0.000) arc (0.00:31.66:0.6);\node[inner sep=0pt,font=\footnotesize] at (1.106,0.314) {$31{,}7^\circ$};\draw[mbgrau,line width=0.5pt] (2.574,1.588) arc (-148.34:-90.00:0.5);\node[inner sep=0pt,font=\footnotesize] at (2.610,1.151) {$\gamma_1$};\node[font=\small,anchor=north east] at (0.000,0.000) {$A$};\node[font=\small,anchor=north west] at (4.450,0.000) {$B$};\node[font=\small,anchor=south] at (3.000,1.850) {$C$};\node[font=\small,anchor=north] at (3.000,0.000) {$D$};\node[font=\small,anchor=west] at (3.050,1.000) {$h_c$};\end{tikzpicture}}{\karo{2}}
\fusshilfe{\mbox{6:~58,3°}}
\end{pfaufg}

% 2018-OS-K4b, herausgelöst (Höhen 1,5 m und 12,0 m nur für 4c, 4d, weggelassen)
\begin{pfaufg}{nach~P10\\’18~\pkt{}~4b}{7.}{}
Lea will die Höhe eines Berges bestimmen. Auf der Bergspitze $C$ steht eine Aussichtsplattform, ihre Spitze heißt $D$. Lea peilt vom Punkt $A$ aus. Der Punkt $B$ liegt auf der Höhe von $A$ senkrecht unter $C$. Die Peillinie zu $C$ ist um $30^\circ$ gegen die Waagerechte geneigt. Zwischen den Peillinien zu $C$ und zu $D$ liegt der Winkel $\gamma = 5^\circ$. Lea sagt: „Der Winkel $\delta$ ist $55^\circ$ groß.“ Begründe, dass das stimmt.
\gzwei{\begin{tikzpicture}[line width=0.7pt,scale=0.8]\draw[dashed] (0, 0)--(4, 0);\draw (4, 0)--(4, 3.2);\draw (0, 0)--(4, 2.0);\draw (0, 0)--(4, 3.2);\draw[mbgrau,line width=0.5pt] (4.000,0.300) arc (90.00:180.00:0.3);\fill (3.908,0.092) circle (0.8pt);\draw[mbgrau,line width=0.5pt] (0.900,0.000) arc (0.00:26.57:0.9);\node[inner sep=0pt,font=\footnotesize] at (1.265,0.299) {$30^\circ$};\draw[mbgrau,line width=0.5pt] (1.968,0.984) arc (26.57:38.66:2.2);\node[inner sep=0pt,font=\footnotesize] at (2.148,1.374) {$\gamma$};\draw[mbgrau,line width=0.5pt] (3.610,2.888) arc (-141.34:-90.00:0.5);\node[inner sep=0pt,font=\footnotesize] at (3.653,2.479) {$\delta$};\node[font=\small,anchor=east] at (0.000,0.000) {$A$};\node[font=\small,anchor=north west] at (4.000,0.000) {$B$};\node[font=\small,anchor=west] at (4.000,2.000) {$C$};\node[font=\small,anchor=south] at (4.000,3.200) {$D$};\end{tikzpicture}}{\karo{4}}
\fusshilfe{\mbox{7:~richtig, $\delta = 55^\circ$}}
\end{pfaufg}

% ===================== Stufe 3 =====================
\Needspace*{0.25\textheight}\pfstufekopf{gleichschenklig: zwei gleiche Winkel}{letzte 5 Jahre: 2023}

% nach Bank winkel-dreiecke-e3-k2-s8-v1 (gleiche Winkel, gleiche Seiten), Durchsicht 11.10.: Ankreuzen mit Skizze (erste der Stufe)
\begin{pfaufg}{}{8.}{}
Kreuze an, welche Seiten gleich lang sind.
\gzwei{\begin{tikzpicture}[line width=0.7pt]\draw (0,0)--(2.4,0)--(1.2,2.46)--cycle;\draw[mbgrau,line width=0.5pt] (0.4,0) arc (0:64:0.4);\node[inner sep=0pt,font=\footnotesize] at (0.66,0.36) {$64^\circ$};\draw[mbgrau,line width=0.5pt] (2.0,0) arc (180:116:0.4);\node[inner sep=0pt,font=\footnotesize] at (1.74,0.36) {$64^\circ$};\node[font=\small,anchor=north east] at (0,0) {$A$};\node[font=\small,anchor=north west] at (2.4,0) {$B$};\node[font=\small,anchor=south] at (1.2,2.46) {$C$};\end{tikzpicture}}{\pfkreuzab{A}{$\overline{AB}$ und $\overline{AC}$}
\pfkreuzab{B}{$\overline{AC}$ und $\overline{BC}$}
\pfkreuzab{C}{$\overline{AB}$ und $\overline{BC}$}}
\fusshilfe{\mbox{8:~B}}
\end{pfaufg}

% 2023-OS-B1g, eigener Wortlaut
\begin{pfaufg}{nach~P10\\’23~\pkt{}~1g}{9.}{}
Im Dreieck $ABC$ sind die Winkel bei $A$ und bei $B$ je $70^\circ$ groß. Gib die Länge der Seite $\overline{BC}$ und die Größe des Winkels $\gamma$ an.
\gzwei{\begin{tikzpicture}[line width=0.7pt]\draw (0, 0)--(2.4, 0)--(1.2,3.297)--cycle;\draw[mbgrau,line width=0.5pt] (0.400,0.000) arc (0.00:70.00:0.4);\node[inner sep=0pt,font=\footnotesize] at (0.614,0.430) {$70^\circ$};\draw[mbgrau,line width=0.5pt] (2.263,0.376) arc (110.00:180.00:0.4);\node[inner sep=0pt,font=\footnotesize] at (1.786,0.430) {$70^\circ$};\draw[mbgrau,line width=0.5pt] (1.029,2.827) arc (-110.00:-70.00:0.5);\node[inner sep=0pt,font=\footnotesize] at (1.200,2.497) {$\gamma$};\node[font=\small,anchor=north east] at (0.000,0.000) {$A$};\node[font=\small,anchor=north west] at (2.400,0.000) {$B$};\node[font=\small,anchor=south] at (1.200,3.297) {$C$};\node[font=\small,anchor=east] at (0.500,1.648) {$6$ cm};\end{tikzpicture}}{\vspace{4mm}\begin{tabular}{@{}r@{\;}l@{}}$\overline{BC}$ & $=$\linie{28mm}\\[5mm] $\gamma$ & $=$\linie{28mm}\end{tabular}}
\fusshilfe{\mbox{9:~6 cm; 40°}}
\end{pfaufg}

% 2023-OS-K7a, herausgelöst (65° bei B und AB = 131,5 cm nur für 7b, weggelassen)
\begin{pfaufg}{nach~P10\\’23~\pkt{}~7a}{10.}{}
Im Dreieck $ABC$ liegt der Punkt $F$ auf der Seite $\overline{AC}$. Die Strecke $\overline{BF}$ steht senkrecht auf $\overline{AC}$. Begründe, dass das graue Dreieck $ABF$ gleichschenklig ist.
\gzwei{\begin{tikzpicture}[line width=0.7pt,scale=0.85]\fill[black!10] (5.0, 0)--(3.0, 2.0)--(3.0, 0)--cycle;\draw (0, 0)--(5.0, 0)--(3.0, 2.0)--cycle;\draw (3.0, 2.0)--(3.0, 0);\draw[mbgrau,line width=0.5pt] (3.300,0.000) arc (0.00:90.00:0.3);\fill (3.092,0.092) circle (0.8pt);\draw[mbgrau,line width=0.5pt] (4.682,0.318) arc (135.00:180.00:0.45);\node[inner sep=0pt,font=\footnotesize] at (4.215,0.325) {$45^\circ$};\draw[mbgrau,line width=0.5pt] (3.000,1.550) arc (-90.00:-45.00:0.45);\node[inner sep=0pt,font=\footnotesize] at (3.325,1.215) {$45^\circ$};\node[font=\small,anchor=north east] at (0.000,0.000) {$C$};\node[font=\small,anchor=north west] at (5.000,0.000) {$A$};\node[font=\small,anchor=south] at (3.000,2.000) {$B$};\node[font=\small,anchor=north] at (3.000,0.000) {$F$};\end{tikzpicture}}{\karo{3}}
\fusshilfe{\mbox{10:~zwei Winkel je 45°}}
\end{pfaufg}

% ===================== Stufe 4 =====================
\Needspace*{0.25\textheight}\pfstufekopf{rechter Winkel: begründen}{letzte 5 Jahre: 2025}

% Bank winkel-dreiecke-e3-k2-s9-v6 (zwei gleichschenklig-rechtwinklige Teildreiecke)
\begin{pfaufg}{}{11.}{}
Die Strecke $\overline{CD}$ steht senkrecht auf $\overline{AB}$. Es gilt $\overline{AD} = \overline{DC} = \overline{DB} = 3$~cm. Begründe, dass der Winkel bei $C$ ein rechter ist.
\gzwei{\begin{tikzpicture}[line width=0.7pt,font=\footnotesize]\draw (0.000,0.000) -- (3.000,0.000) -- (1.500,1.500) -- cycle;\draw (1.500,1.500) -- (1.500,0.000);\draw[mbgrau,line width=0.5pt] (1.500,0.300) arc (90.00:180.00:0.3);\fill (1.405,0.095) circle (0.8pt);\node[inner sep=0pt] at (-0.231,-0.096) {$A$};\node[inner sep=0pt] at (3.231,-0.096) {$B$};\node[inner sep=0pt] at (1.500,1.750) {$C$};\node[inner sep=0pt] at (1.500,-0.220) {$D$};\end{tikzpicture}}{\karo{3}}
\fusshilfe{\mbox{11:~45° + 45° = 90°}}
\end{pfaufg}

% 2025-OS-K2c, herausgelöst (Sternchenaufgabe; Symmetrieachsen 2a und Fläche 2b weggelassen)
\begin{pfaufg}{nach~P10\\’25~\pkt{}~2c\,$\star$}{12.}{}
Die Figur zeigt das Drachenviereck $ABCD$. Seine Diagonalen $\overline{AC}$ und $\overline{BD}$ sind $50$~cm und $80$~cm lang und schneiden sich im Punkt $F$. Die Strecke $\overline{DF}$ ist $25$~cm lang. Begründe, dass $\gamma = 90^\circ$ ist.
\gzwei{\begin{tikzpicture}[line width=0.7pt,scale=0.75]\draw (-1.25, 0)--(0, -2.75)--(1.25, 0)--(0, 1.25)--cycle;\draw (-1.25, 0)--(1.25, 0);\draw (0, -2.75)--(0, 1.25);\draw[mbgrau,line width=0.5pt] (0.220,0.000) arc (0.00:90.00:0.22);\fill (0.092,0.092) circle (0.8pt);\draw[mbgrau,line width=0.5pt] (-0.247,1.003) arc (-135.00:-45.00:0.35);\node[font=\footnotesize] at (-0.2,0.7) {$\gamma$};\node[font=\small,anchor=east] at (-1.250,0.000) {$A$};\node[font=\small,anchor=north] at (0.000,-2.750) {$B$};\node[font=\small,anchor=west] at (1.250,0.000) {$C$};\node[font=\small,anchor=south] at (0.000,1.250) {$D$};\node[font=\small,anchor=north west] at (0.000,0.000) {$F$};\end{tikzpicture}}{\karo{4}}
\fusshilfe{\mbox{12:~45° + 45° = 90°}}
\end{pfaufg}

\pfblattende{11}{26 FOR & 25 & 24 & 23 & 22 & 21 & 20 & 19 & 18 & 17 & 16}%
{1c 1i & 2c & & 1g 2a 7a & 5c & 1i & & & 4b & 4a & 1e}%
{Weiter: Sinussatz – in jedem Dreieck}
\end{document}
